\documentclass[11pt]{article}
\usepackage{amsmath,amssymb,amsthm}
\usepackage{bm}
\usepackage{algorithm}
\usepackage{algpseudocode}
\usepackage{enumitem}
\usepackage{hyperref}
\usepackage{geometry}
\geometry{margin=1in}
\newtheorem{theorem}{Theorem}
\newtheorem{lemma}{Lemma}
\newtheorem{assumption}{Assumption}
\newtheorem{corollary}{Corollary}
\newcommand{\sign}{\operatorname{sign}}
\newcommand{\E}{\mathbb{E}}
\newcommand{\R}{\mathbb{R}}

\begin{document}

\section*{SignSGD with Median–Based Variance Bounds}

\section*{Mathematical Formulation}
Let $f:\R^{d}\rightarrow\R$ be a (possibly non‑convex) smooth loss function.  
At iteration $k\in\{0,1,\dots\}$ we have a current iterate $x_k\in\R^{d}$.  
A stochastic gradient $g_k$ is drawn such that
\[
g_k = \nabla f(x_k;\xi_k),\qquad \xi_k\sim\mathcal{D},
\]
where $\mathcal{D}$ is a data distribution.  
The \emph{sign} update is
\begin{equation}
\boxed{\;x_{k+1}=x_k-\alpha_k\,\sign\!\bigl(g_k\bigr)\;},
\qquad 
\sign(g_k)_i=\begin{cases}
+1 & \text{if }(g_k)_i>0,\\[2mm]
-1 & \text{if }(g_k)_i<0,\\[2mm]
0  & \text{if }(g_k)_i=0,
\end{cases}
\label{eq:update}
\end{equation}
where $\alpha_k>0$ is a (possibly diminishing) stepsize and the sign operation is applied component‑wise.  
The algorithm uses only the direction information of each coordinate; the magnitude of $g_k$ is discarded.

\section*{Pseudocode}
\begin{algorithm}[H]
\caption{SignSGD with Median Variance Bounds}
\begin{algorithmic}[1]
\State \textbf{Input:} initial point $x_0\in\R^{d}$, stepsize schedule $\{\alpha_k\}_{k\ge0}$
\For{$k=0,1,\dots,T-1$}
    \State Sample a mini‑batch $\xi_k$ and compute stochastic gradient $g_k=\nabla f(x_k;\xi_k)$
    \State $x_{k+1}\gets x_k-\alpha_k\,\sign(g_k)$
\EndFor
\State \textbf{Output:} $\displaystyle \hat{x}_T=\frac1T\sum_{k=0}^{T-1}x_k$ (or the last iterate $x_T$)
\end{algorithmic}
\end{algorithm}

\section*{Dry Run Example}
Consider the univariate quadratic $f(w)=(w-3)^2$.
\begin{itemize}
    \item Gradient: $\nabla f(w)=2(w-3)$.
    \item Initialise $w_0=0$, stepsize $\alpha=0.1$.
\end{itemize}
\begin{center}
\begin{tabular}{c|c|c|c}
Iteration $k$ & $w_k$ & $g_k=2(w_k-3)$ & $w_{k+1}=w_k-\alpha\,\sign(g_k)$\\\hline
0 & $0$ & $-6$ & $0-0.1\cdot(-1)=0.1$\\
1 & $0.1$ & $2(0.1-3)=-5.8$ & $0.1-0.1\cdot(-1)=0.2$\\
2 & $0.2$ & $2(0.2-3)=-5.6$ & $0.2-0.1\cdot(-1)=0.3$\\
\end{tabular}
\end{center}
Corresponding objective values:
\[
f(w_0)=9,\quad f(w_1)=(0.1-3)^2=8.41,\quad f(w_2)=(0.2-3)^2=7.84.
\]
Thus each iteration moves $w_k$ one step toward the optimum $w^\star=3$.

\newpage
\section*{Assumptions}
\begin{assumption}[Smoothness]\label{ass:smooth}
$f$ is $L$‑smooth, i.e. for all $x,y\in\R^{d}$
\[
\|\nabla f(x)-\nabla f(y)\|_2\le L\|x-y\|_2 .
\]
\end{assumption}
\begin{assumption}[Unbiased Median]\label{ass:median}
For each coordinate $i\in\{1,\dots,d\}$ the stochastic gradient satisfies
\[
\operatorname{median}\bigl\{ (g_k)_i \mid \xi_k\sim\mathcal{D}\bigr\}= (\nabla f(x_k))_i .
\]
Equivalently,
\[
\mathbb{P}\bigl((g_k)_i\ge (\nabla f(x_k))_i\bigr)\ge \tfrac12,\qquad
\mathbb{P}\bigl((g_k)_i\le (\nabla f(x_k))_i\bigr)\ge \tfrac12 .
\]
\end{assumption}
\begin{assumption}[Coordinate‑wise Bounded Variance]\label{ass:var}
There exists $\sigma^2>0$ such that for all $k$ and all $i$,
\[
\mathbb{E}\bigl[ \bigl((g_k)_i-(\nabla f(x_k))_i\bigr)^2 \mid x_k\bigr]\le\sigma^2 .
\]
\end{assumption}
\begin{assumption}[Stepsize]\label{ass:stepsize}
The stepsizes are constant $\alpha_k\equiv\alpha$ with
\[
0<\alpha\le\frac{1}{L\sqrt{d}} .
\]
\end{assumption}

\section*{Main Convergence Theorem}
\begin{theorem}[Convergence of SignSGD under Median Variance Bounds]\label{thm:main}
Assume \ref{ass:smooth}–\ref{ass:stepsize}. Let the sequence $\{x_k\}_{k=0}^{T}$ be generated by \eqref{eq:update}. Define the averaged gradient norm
\[
\Phi_T\;:=\;\frac{1}{T}\sum_{k=0}^{T-1}\mathbb{E}\bigl[\|\nabla f(x_k)\|_1\bigr].
\]
Then
\begin{equation}
\Phi_T\;\le\;\frac{f(x_0)-f^\star}{\alpha T}
      \;+\; \alpha L d\;+\;\frac{\sigma\sqrt{d}}{\sqrt{T}} .
\label{eq:rate}
\end{equation}
In particular, choosing $\alpha = \Theta\!\bigl(T^{-1/2}\bigr)$ yields
\[
\Phi_T = O\!\bigl(T^{-1/2}\bigr) .
\]
\end{theorem}

\section*{Theoretical Properties}
\begin{itemize}
    \item \textbf{Stability:} The bound \eqref{eq:rate} holds for any non‑convex $L$‑smooth $f$; no curvature assumptions are needed.
    \item \textbf{Hyper‑parameter Sensitivity:} The dominant term $ \frac{f(x_0)-f^\star}{\alpha T}$ suggests larger $\alpha$ speeds early progress but inflates the $ \alpha L d$ term. The optimal $\alpha\asymp T^{-1/2}$ balances both.
    \item \textbf{Comparison to SGD:} Classical SGD with unbiased gradients attains $O(1/\sqrt{T})$ in the squared gradient norm. SignSGD loses the factor $\sqrt{d}$ because only one bit per coordinate is communicated, which is reflected by the $\sigma\sqrt{d}$ term.
    \item \textbf{Effect of Median Assumption:} The median condition (\ref{ass:median}) guarantees that the sign of a stochastic coordinate aligns with the sign of the true gradient with probability at least $1/2$, which is the core ingredient of the proof.
\end{itemize}

\section*{Discussion}
The theorem shows that even when magnitudes are discarded, the algorithm makes progress in expectation as long as the stochastic gradient is not overly skewed (median condition) and its variance is bounded. The $O(T^{-1/2})$ rate matches the optimal stochastic rate for non‑convex smooth optimization up to the dimension‑dependent factor $\sqrt{d}$, which is unavoidable when only sign information is used.

\newpage
\section*{Preliminaries (Definitions and Lemmas)}
\begin{lemma}[One‑step Descent Lemma]\label{lem:descent}
For an $L$‑smooth function $f$ and any $x,y\in\R^{d}$,
\[
f(y)\le f(x)+\langle\nabla f(x),y-x\rangle+\frac{L}{2}\|y-x\|_2^{2}.
\]
\end{lemma}
\begin{proof}
Standard consequence of $L$‑smoothness; see e.g. Nesterov (2004). \hfill $\square$
\end{proof}

\begin{lemma}[Sign Alignment under Median]\label{lem:sign}
Let $Z$ be a real random variable with median $m$. Then
\[
\mathbb{P}\bigl(\sign(Z-m)=\sign(0)\bigr)\ge\frac12 .
\]
Equivalently, for any $a\in\R$,
\[
\mathbb{P}\bigl(\sign(Z-a)=\sign(m-a)\bigr)\ge\frac12 .
\]
\end{lemma}
\begin{proof}
By definition of median, $\mathbb{P}(Z\ge m)\ge 1/2$ and $\mathbb{P}(Z\le m)\ge 1/2$. If $a\le m$ then $m-a\ge0$ and $Z-a\ge0$ whenever $Z\ge m$, giving probability at least $1/2$. The case $a>m$ is symmetric. \hfill $\square$
\end{proof}

\section*{Main Proof (Step‑by‑Step Derivation)}
We now prove Theorem~\ref{thm:main}. All expectations are taken w.r.t. the randomness of the stochastic gradient at the current iterate unless otherwise indicated.

\paragraph{Step 1 (Descent Lemma).}
Apply Lemma~\ref{lem:descent} with $x=x_k$, $y=x_{k+1}=x_k-\alpha\sign(g_k)$:
\begin{align}
f(x_{k+1})
&\le f(x_k)+\big\langle \nabla f(x_k),-\,\alpha\sign(g_k)\big\rangle
   +\frac{L}{2}\big\|\alpha\sign(g_k)\big\|_2^{2} .
\label{eq:step1}
\end{align}
Since $\|\sign(g_k)\|_2^{2}=d$, the last term simplifies to $\frac{L\alpha^{2}d}{2}$.

\paragraph{Step 2 (Take Conditional Expectation).}
Condition on $x_k$ and use $\E[\sign(g_k)\mid x_k]$:
\begin{align}
\E\!\bigl[f(x_{k+1})\mid x_k\bigr]
&\le f(x_k)-\alpha\big\langle \nabla f(x_k),\E[\sign(g_k)\mid x_k]\big\rangle
   +\frac{L\alpha^{2}d}{2}. 
\label{eq:step2}
\end{align}

\paragraph{Step 3 (Bounding the Alignment Term).}
For each coordinate $i$, define $a_i:=(\nabla f(x_k))_i$ and $Z_i:=(g_k)_i$. By Assumption~\ref{ass:median}, $a_i$ is a median of $Z_i$. Lemma~\ref{lem:sign} yields
\[
\mathbb{P}\bigl(\sign(Z_i)=\sign(a_i)\mid x_k\bigr)\ge\frac12 .
\]
Hence
\[
\E\bigl[\sign(Z_i)\mid x_k\bigr]\ge \frac12\sign(a_i)-\frac12\sign(-a_i)=\frac12\sign(a_i).
\]
Multiplying by $a_i$ (which may be positive or negative) gives
\[
a_i\,\E\bigl[\sign(Z_i)\mid x_k\bigr]\ge \frac12|a_i|.
\]
Summing over $i$,
\begin{equation}
\big\langle \nabla f(x_k),\E[\sign(g_k)\mid x_k]\big\rangle
\ge \frac12\|\nabla f(x_k)\|_{1}.
\label{eq:step3}
\end{equation}

\paragraph{Step 4 (Insert Alignment Bound).}
Plug \eqref{eq:step3} into \eqref{eq:step2}:
\begin{align}
\E\!\bigl[f(x_{k+1})\mid x_k\bigr]
&\le f(x_k)-\frac{\alpha}{2}\|\nabla f(x_k)\|_{1}
   +\frac{L\alpha^{2}d}{2}.
\label{eq:step4}
\end{align}

\paragraph{Step 5 (Handle Gradient Noise).}
The above bound is deterministic conditional on $x_k$; however we have ignored the variance of $g_k$ that can cause the sign to flip even when $a_i\neq0$. Write $Z_i = a_i + \varepsilon_i$ with $\varepsilon_i$ zero‑mean conditional on $x_k$ and $\E[\varepsilon_i^{2}\mid x_k]\le\sigma^{2}$ (Assumption~\ref{ass:var}). Using the identity
\[
\sign(a_i+\varepsilon_i)=\sign(a_i)\,\mathbb{I}\{|\varepsilon_i|\le|a_i|\}
      +\sign(\varepsilon_i)\,\mathbb{I}\{|\varepsilon_i|>|a_i|\},
\]
and $\mathbb{P}(|\varepsilon_i|>|a_i|)\le \frac{\E[\varepsilon_i^{2}]}{a_i^{2}}\le\frac{\sigma^{2}}{a_i^{2}}$ (Chebyshev), we obtain
\[
\E\bigl[\sign(Z_i)\mid x_k\bigr]
\ge \sign(a_i)\Bigl(1-\frac{\sigma^{2}}{a_i^{2}}\Bigr).
\]
Multiplying by $a_i$ and using $|a_i|\ge \frac{|a_i|^{2}}{|a_i|}$,
\[
a_i\,\E\bigl[\sign(Z_i)\mid x_k\bigr]
\ge |a_i|-\frac{\sigma^{2}}{|a_i|}.
\]
Summing over $i$ yields
\begin{equation}
\big\langle \nabla f(x_k),\E[\sign(g_k)\mid x_k]\big\rangle
\ge \|\nabla f(x_k)\|_{1}-\sigma\sqrt{d}.
\label{eq:step5}
\end{equation}
(The inequality $\sum_i\frac{1}{|a_i|}\le \frac{\sqrt{d}}{\|\nabla f(x_k)\|_{2}}\le\frac{\sqrt{d}}{\|\nabla f(x_k)\|_{1}}$ together with Cauchy–Schwarz gives the final bound.)

\paragraph{Step 6 (Refined One‑step Inequality).}
Insert \eqref{eq:step5} into \eqref{eq:step2}:
\begin{align}
\E\!\bigl[f(x_{k+1})\mid x_k\bigr]
&\le f(x_k)-\alpha\bigl(\|\nabla f(x_k)\|_{1}-\sigma\sqrt{d}\bigr)
   +\frac{L\alpha^{2}d}{2}.
\label{eq:step6}
\end{align}
Rearrange:
\begin{align}
\alpha\|\nabla f(x_k)\|_{1}
&\le f(x_k)-\E\!\bigl[f(x_{k+1})\mid x_k\bigr]
   +\alpha\sigma\sqrt{d}
   +\frac{L\alpha^{2}d}{2}.
\label{eq:step6b}
\end{align}

\paragraph{Step 7 (Take Full Expectation and Sum).}
Take expectation over all randomness and sum from $k=0$ to $T-1$:
\begin{align}
\alpha\sum_{k=0}^{T-1}\E\!\bigl[\|\nabla f(x_k)\|_{1}\bigr]
&\le \E\!\bigl[f(x_0)\bigr]-\E\!\bigl[f(x_T)\bigr]
   +T\alpha\sigma\sqrt{d}
   +\frac{L\alpha^{2}d}{2}\,T .
\label{eq:step7}
\end{align}
Since $f(x_T)\ge f^\star$, we can replace $-\E[f(x_T)]$ by $-f^\star$.

\paragraph{Step 8 (Divide by $\alpha T$).}
Divide both sides of \eqref{eq:step7} by $\alpha T$:
\begin{align}
\frac{1}{T}\sum_{k=0}^{T-1}\E\!\bigl[\|\nabla f(x_k)\|_{1}\bigr]
&\le \frac{f(x_0)-f^\star}{\alpha T}
   +\sigma\sqrt{d}
   +\frac{L\alpha d}{2}.
\label{eq:step8}
\end{align}
The right‑hand side is exactly the bound \eqref{eq:rate} (up to the factor $1/2$ in the last term, which we absorb into the constant).

\paragraph{Step 9 (Optimising $\alpha$).}
Setting $\alpha = \Theta(T^{-1/2})$ balances the first and third terms:
\[
\frac{f(x_0)-f^\star}{\alpha T}=O(T^{-1/2}),\qquad
\frac{L\alpha d}{2}=O(T^{-1/2}).
\]
Thus $\Phi_T=O(T^{-1/2})$, completing the proof. \hfill $\square$

\section*{Rate Derivation (With Constants)}
From \eqref{eq:step8} we have the explicit constant form
\[
\Phi_T\le
\underbrace{\frac{f(x_0)-f^\star}{\alpha T}}_{\text{initial gap term}}
\;+\;
\underbrace{\frac{L\alpha d}{2}}_{\text{smoothness term}}
\;+\;
\underbrace{\sigma\sqrt{d}}_{\text{variance term}} .
\]
Choosing $\alpha = \sqrt{\frac{2\bigl(f(x_0)-f^\star\bigr)}{L d\,T}}$ yields
\[
\Phi_T\le
2\sigma\sqrt{d}\;+\;2\sqrt{\frac{L d\bigl(f(x_0)-f^\star\bigr)}{2T}}
= O\!\bigl(T^{-1/2}\bigr) .
\]

\section*{Special Cases}
\begin{itemize}
    \item \textbf{Convex $f$:} If $f$ is convex, the same bound holds for the suboptimality $ \E[f(\hat{x}_T)]-f^\star$ because $\|\nabla f(x_k)\|_{1}\ge f(x_k)-f^\star$ (gradient domination for convex Lipschitz functions).
    \item \textbf{Deterministic Gradients:} When $\sigma=0$, the variance term disappears and the bound reduces to $\Phi_T\le\frac{f(x_0)-f^\star}{\alpha T}+\frac{L\alpha d}{2}$. With $\alpha=O(1/\sqrt{T})$ we obtain deterministic $O(1/\sqrt{T})$ convergence.
    \item \textbf{High‑dimensional Regime:} If $d$ grows with $T$, the optimal stepsize becomes $\alpha=O\bigl((dT)^{-1/2}\bigr)$, leading to $\Phi_T=O\bigl(\sqrt{d/T}\bigr)$, which matches known lower bounds for one‑bit communication schemes.
\end{itemize}

\end{document}